% !TEX root = ../main.tex

\chapter{The Thesis in Brief and What It Left Open}
\label{ch:thesis-recap}

\section{The Thesis in Brief}
\label{sec:thesis-brief}

A robot acting in a cluttered scene often cannot see everything that matters: a
target object may be hidden behind others. The thesis \cite{alassiri2026thesis}
extends classical task and motion planning to this case by a \emph{Boxel}
representation. A Boxel is a box-shaped cell of the workspace. Each detected
object gets its own cell, each region that an object hides from the camera (a
\emph{shadow}) is a cell of its own, and the remaining free space is covered with
a few coarse cells. The robot's belief is a set of yes-or-no facts, one per cell,
that say whether the cell's content is known (Know-If fluents). A classical
planner, PDDLStream \cite{garrett2018pddlstream} with Fast Downward
\cite{helmert2006fast}, reasons over these facts and plans sensing in the same
search as the manipulation. The system plans optimistically, assuming that a
hidden target is in the cell it is about to look at, acts, senses, and replans
when an assumption fails. It runs on a simulated Franka Panda in PyBullet
\cite{coumans2021pybullet}.

The thesis measured the representation by ablation: the same planner on the
adaptive partition against a uniform grid at the same cell scale. Its headline
table, reproduced in \Cref{tab:thesis-headline} with four rows recounted, shows that the adaptive
partition solved more scenes in every goal, most sharply on finding a hidden
cube and stacking it on a tray.

\begin{table}[tbp]
\centering
\small
\begin{tabular}{llrr}
\toprule
goal & variant & success rate & mean plan time (s) \\
\midrule
find-and-tray-stack & semantic         & 62 of 90, 68.9\,\% &  99.99 \\
find-and-tray-stack & semantic+mbs0.05 & 58 of 90, 64.4\,\% &  79.42 \\
find-and-tray-stack & uniform          & 27 of 90, 30.0\,\% & 705.07 \\
\midrule
holding             & semantic         & 65.6\,\% & 134.30 \\
holding             & semantic+mbs0.05 & 66.7\,\% &  56.64 \\
holding             & uniform          & 40 of 90, 44.4\,\% & 401.32 \\
\midrule
stack               & semantic         & 64.4\,\% &  17.16 \\
stack               & semantic+mbs0.05 & 64.4\,\% &   3.10 \\
stack               & uniform          & 35.6\,\% & 924.13 \\
\bottomrule
\end{tabular}
\caption{The thesis's headline table \cite{alassiri2026thesis}: success rate and
mean success-only planning time per goal and variant, 90 episodes each (30 seeds
and three difficulty levels). Four rows are recounted on the first attempt of each
declared cell (\ID{recount}{recount of 14 November}): where the declared scene failed to
build, a re-rolled scene ran in its place (5, 7 and 4 cells of find-and-tray-stack, all
successes), and two runs of find-and-tray-stack semantic and four of holding uniform were
cut by our runner's restart; these count as failures. The thesis printed 75.6, 72.2 and
34.4\,\% and 47.8\,\%. The other rows are as printed.}
\label{tab:thesis-headline}
\end{table}

\section{The Thesis's Comparison with TAMPURA}
\label{sec:thesis-tampura}

The thesis compared its \texttt{holding} task with the hidden-object task of
TAMPURA \cite{curtis2024partially}, which the paper calls Partial Observability
(task C) and the released code calls \fdice. It re-ran TAMPURA's released
environment on the thesis machine over 20 seeds and reported 11 of 20
successes (\PW{4}), a mean of 166 seconds per successful episode against 144.4
seconds for the thesis system's wall clock, and a success rate of 65.6\,\% for the
thesis system against the at least 63\,\% that TAMPURA's published return of 0.63
implies
\cite{alassiri2026thesis,curtis2024partially}. The thesis itself calls this
comparison architectural: the two systems ran on different environments, so it
benchmarks neither of them.

\section{What the Thesis Said About Its Own Limits}
\label{sec:thesis-limits}

The thesis's limitations section states the simplifications behind its numbers
\cite{alassiri2026thesis}. Perception was an oracle that read exact poses from
the simulator, and the camera was fixed. Grasping was constraint-based: a held
object was rigidly welded to the gripper, so it could not slip or be dropped,
and only top-down grasps at one fixed height were sampled. A short hard-coded
lift followed every place or stack. Collisions during execution were logged but
did not trigger a replan. The shadow cells under-covered the region an object
hides, which the thesis names as the dominant cause of the search-exhausted
failures (13 of 18 cells). The scenes had 6 to 9 cm cubes, and the TAMPURA task
was the closest analogue, not the same problem.

\section{What Was Asked Next}
\label{sec:asked-next}

After the thesis, the supervisors' list of missing pieces became ten workstreams
in the audit file. The wording of the first is the supervisors' as the audit
records it: ``no workarounds and cuts any more. The ONLY thing that may stay is
the hardcoded lift after a pick or place action'' (\PW{1}). The paper that the
work feeds has a submission date of 4 December 2026 (\BF{1}).
\Cref{tab:asked-next} lists the workstreams and where this report treats them.

\begin{table}[tbp]
\centering
\small
\begin{tabularx}{\textwidth}{@{}lYYl@{}}
\toprule
workstream & what was asked & what happened & chapter \\
\midrule
\PW{1} & remove the weld and every oracle read; build real perception and a real pose sampler & landed in steps from 20 August to 19 September 2026, with 32 field findings (F1 to F32) & \ref{ch:system} \\
\PW{2} & prove the stacking drop is execution, detect physical failure, fail honestly & monitor landed 19 September; root cause for height 4 found 20 September & \ref{ch:system} \\
\PW{3} & all objects must spawn within reachable range & landed 19 September; reach measured & \ref{ch:system} \\
\PW{4} & run our pipeline inside TAMPURA's environment for a valid comparison & restarted 25 October on a new branch; \BF{1} to \BF{68} & \ref{ch:comparison} \\
\PW{5} & consider adopting TAMPURA's controllers & open, a decision (marked T2 in the audit) on adopting their grasp library; the campaign of \PW{4} leaves it out of scope & none \\
\PW{6} & a real-robot experiment & a remark: not needed for ICAPS, nice to have & none \\
\PW{7} & related work 2024 to 2026 & open & \ref{ch:open-questions} \\
\PW{8} & comparable results from other approaches in the same environment & open; TAMPURA and LW1 are the candidates already on the list & \ref{ch:open-questions} \\
\PW{9} & full partially observable planning with LW1 and PDDLStream, 3-D Wumpus & stretch goal; the entry names a Wumpus-World baseline as its first rung, and paths 4 and 10 work in the Wumpus world & \ref{ch:progress} \\
\PW{10} & division of labour and framing of the paper & a remark; framing question open & \ref{ch:open-questions} \\
\bottomrule
\end{tabularx}
\caption{The ten workstreams of the audit file.}
\label{tab:asked-next}
\end{table}
